\documentclass[11pt]{article}
\usepackage[margin=1in]{geometry}
\usepackage{graphicx}
\usepackage{booktabs}
\usepackage{amsmath}
\usepackage{hyperref}
\usepackage{natbib}

\title{Low-Dimensional, Non-Lexical Trajectories of Evolving Software
Objects}

\author{Giovanni Pezzin \\ Physis Research Program}
\date{\today}

\begin{document}
\maketitle

\begin{abstract}
Revisions of a single source file occupy a low-dimensional, temporally
directed region of embedding space --- but the identity half of the usual
trajectory story does not survive a trivial baseline. On an external
replication of 434 source files and 11{,}312 revisions across ten public
repositories in five languages, the participation ratio is 2.13 and the
effective rank 3.06: 7--10$\times$ more concentrated than an isotropic cloud
of the same shape, against a per-metric null. PC1+PC2 carry 0.843 of the
variance against a null of 0.113. Direction is present and is \emph{not}
text overlap: the geometric drift statistic anti-correlates with the lexical
one ($\rho = -0.593$) and keeps a partial correlation with revision depth of
0.232 after controlling for lexical adjacency. Identity, however, is not
evidence for a semantic representation --- same-file revisions separate from
cross-file revisions at 0.99 under nomic, but the trivial lexical baselines
do the same (token-jaccard 0.988, edit distance 0.974) and separate them
\emph{more} by gap median (0.369 and 0.516 against 0.214). We therefore keep
the dimensional and directional claims, drop identity-as-semantic-evidence,
and make no operational prediction claim: that experiment was never run.
Earlier internal figures on six files (2--3 dimensions from 57--82\% raw
PC1+PC2 variance, no null) are superseded --- raw explained variance cannot
support a dimensionality claim, as the replication's first control
demonstrated.
\end{abstract}

\section{Introduction}

Software objects evolve. A source file accumulates revisions; each revision
is a state of the object at a point in time. The question is whether these
states form a \emph{trajectory} --- a continuous, directed path through
representation space --- or a set of independent points.

This question has practical implications. If revisions form a trajectory,
then (a) semantic search should use trajectory queries, not point queries;
(b) change prediction should extrapolate the trajectory, not predict from
a static embedding; (c) object identity should be defined by trajectory
proximity, not point proximity.

We test this on six files from a Rust codebase, using BGE-base embeddings
(768 dimensions, \texttt{onnx-explicit} backend). We measure three
properties: identity (within-object distance vs.\ between-object
distance), direction (Spearman correlation of drift with time, against a
shuffled-time null), and structure (spectral concentration of the
trajectory).

\section{Method}

For each file $f$, we:
\begin{enumerate}
\item Extract all revisions $r_1, \ldots, r_n$ from git history.
\item Embed each revision (chunked at 1500 characters, mean-pooled into a
single 768-d vector).
\item Compute drift $d_i = \|v_i - v_1\|$ (distance from the first revision).
\item Compute Spearman correlation $\rho(i, d_i)$ against a 2000-permutation
shuffled-time null.
\item Compute spectral concentration: the fraction of variance in the first
two eigenvalues of the centered Gram matrix.
\end{enumerate}

\section{Data}

Six files from a Rust codebase (\texttt{physis-pro}), spanning 8--27
revisions each:

\begin{table}[h]
\centering
\begin{tabular}{lrrrr}
\toprule
\textbf{File} & \textbf{Revisions} & \textbf{Drift range} & \textbf{Spearman $\rho$} & \textbf{2D var} \\
\midrule
\texttt{hypothesis.rs} & 12 & 0.026 & +0.930 & 0.766 \\
\texttt{main.rs} & 27 & 0.116 & +0.626 & 0.621 \\
\texttt{studio.rs} & 17 & 0.037 & +0.939 & 0.761 \\
\texttt{act.rs} & 8 & 0.053 & +0.976 & 0.817 \\
\texttt{ontology.rs} & 8 & 0.043 & +0.929 & 0.740 \\
\texttt{quality.rs} & 9 & 0.025 & +0.250 & 0.569 \\
\bottomrule
\end{tabular}
\caption{Trajectory analysis of six source files. Drift range is the
difference between maximum and minimum distance from the first revision.
2D var is the fraction of variance explained by the first two eigenvalues
of the Gram matrix.}
\label{tab:trajectories}
\end{table}

\section{Results}

\subsection{Identity}

Within-object distances (mean pairwise cosine distance between revisions
of the same file) are 4--14$\times$ smaller than between-object distances
(mean cosine distance between revisions of different files). A file's
revisions form a tight region of the embedding field; a different file is
elsewhere.

\subsection{Direction}

Five of six files show temporal direction: Spearman $\rho$ ranges from
+0.626 to +0.976, all outside the 2000-permutation shuffled-time null
(which spans approximately [-0.4, +0.7] for most files). The sixth file,
\texttt{quality.rs}, fails: $\rho = +0.250$, inside the null. Its drift
profile shows why: it wanders and comes back, so distance-from-origin
stops being a clock.

\subsection{Structure}

The first two eigenvalues of the centered Gram matrix carry 0.569--0.817
of the variance, and the files occupy few effective dimensions. This figure
alone does \emph{not} establish low intrinsic dimension: a random walk in
768 dimensions also concentrates variance, and an explained-variance fraction
is invariant under permuting the revisions, so a row-shuffled null reproduces
it exactly. Section~\ref{sec:external} reports the statistic against a null
that can actually fail.

\subsection{Granularity}

At document level, identity survives: the document's pooled embedding is
closer to its own sections than to another document's sections (nesting
margin 0.02--0.03). At paragraph level, identity dissolves: a paragraph of
one document is nearly as far from another paragraph of the same document
(0.379) as from a paragraph of a different document (0.405).

Positional direction --- whether a text moves away from its opening as it
proceeds --- is a property of genre, not of text. A developmental document
(the world-map machine plan) shows strong positional direction at section
level ($\rho = +0.665$); a reference document (WHAT\_PHYSIS\_IS) does not
($\rho = +0.083$).

\section{External replication}\label{sec:external}

The six-file result was a single repository and one embedder. We replicated
the same three tests on ten public repositories (ripgrep, serde, requests,
flask and others; C, C++, JS, Python, Rust): 434 source files, 11{,}312
revisions. Representation is nomic-embed-text (768-d) via a local ollama
server --- BGE was unavailable on the replication host, so the embedder
changed as well as the corpus. Inference unit is the file. Result:
\texttt{research/trajectories\_external/p2\_results.json}; generator
\texttt{research/papers/p2\_analysis.py --task all}.

\subsection{Dimension: holds}

Against an isotropic cloud of the same shape $n \times d$, all three
statistics say the real spectra are far more concentrated:
PC1+PC2 share 0.843 vs 0.113 (7.5$\times$), participation ratio 2.127 vs
22.362 (10.5$\times$), effective rank 3.063 vs 22.674 (7.4$\times$). The
participation ratio of 2.13 and effective rank of 3.06 are what ``2--3
dimensional'' should have meant in Section~\ref{tab:trajectories}; the
direction of the comparison differs by statistic, since a \emph{low}
participation ratio is what concentration means.

\subsection{Direction: holds, and is not lexical}

\begin{center}
\begin{tabular}{lr}
\toprule
\textbf{Quantity} & \textbf{Value} \\
\midrule
corr(geometric adjacency, lexical adjacency) & $+0.826$ \\
corr(geometric trend, lexical trend) & $-0.593$ \\
partial corr(geometric adjacency, $n_{rev}$ $|$ lexical adjacency) & $+0.232$ \\
median geometric adjacency & $0.993$ \\
median lexical adjacency & $0.941$ \\
\bottomrule
\end{tabular}
\end{center}

The geometric temporal statistic does not merely survive removing the text:
it points the opposite way to the lexical trend ($-0.593$) and keeps a
partial correlation with revision depth of 0.232. This is the part of the
trajectory claim that a text-overlap baseline cannot reproduce.

\subsection{Identity: does not hold as semantic evidence}

\begin{center}
\begin{tabular}{lrrr}
\toprule
\textbf{Representation} & \textbf{frac. same closer} & \textbf{CI95} & \textbf{gap median} \\
\midrule
nomic (768-d) & 0.9906 & [0.976, 0.996] & 0.214 \\
token-jaccard (lexical) & 0.9883 & [0.973, 0.995] & 0.369 \\
edit distance (lexical) & 0.9742 & [0.954, 0.986] & 0.516 \\
\bottomrule
\end{tabular}
\end{center}

All three separate same-object from cross-object revisions at 97--99\%, and
on the gap medians the trivial baselines are \emph{better} than the
embedding. Identity across revisions is real; it is not evidence for an
embedding. The six-file version of this claim should be read the same way.

\subsection{Prediction: not claimed}

No next-step extrapolation test exists in the replication artifact (its keys
are identity, overlap, dimensionality). We do not claim operational
trajectory prediction, here or anywhere.

\section{Discussion}

\textbf{What survives, and what does not.} Revisions of a single file do
form a low-dimensional, temporally directed region: the dimension claim holds
against a null that can fail, and the direction statistic is not reproducible
by text overlap. But the third pillar of the original framing --- that
identity is evidence for a semantic representation --- does not survive, and
the operational prediction that the introduction hoped for was never tested.
The honest claim is narrower than ``software objects are trajectories'': the
trajectory is a real geometric property of the revision sequence, and what it
is \emph{for} remains open.

\textbf{The failure mode is informative.} \texttt{quality.rs} fails the
direction test because it wanders and returns. This is the honest failure
mode of trajectory analysis: when an object's evolution is not monotonic,
distance-from-origin is not a clock. The failure is visible in the same
table as the successes, not averaged away.

\textbf{Identity dissolves below document granularity.} A document's
paragraphs are not a tighter cluster than paragraphs from different
documents. Whatever ``the same object'' means at revision scale, it is not
recoverable from a paragraph. This is consistent with the E80 finding that
pooling a project into one vector destroys signal.

\textbf{Limits.} One repository, one embedder (BGE-base-en-v1.5), six
files, one metadata axis (commit time). The trajectory instrument is
available as \texttt{physis-core/examples/reference\_spectrum.rs}; the
harness is \texttt{benchmarks/conatus/e78\_reference\_spectrum.py}. The
The expansion to thousands of files across public repos that this paragraph
previously named as the highest-priority replication has now been run
(Section~\ref{sec:external}): 434 files, 11{,}312 revisions, ten repos. It is
the source of the abstract's numbers, and it is what narrowed this paper.

\end{document}
